% Modos de conexión TeleCon — qué significa el ajuste del teléfono.
% Compilar: pdflatex -output-directory=../build telecon_connection_modes.tex
% PDF entregado: ../build/connection_modes_esp.pdf
\documentclass[12pt,a4paper]{article}

\usepackage[margin=2.4cm]{geometry}
\usepackage[T1]{fontenc}
\usepackage[utf8]{inputenc}
\usepackage[spanish]{babel}
\usepackage{lmodern}
\usepackage{booktabs}
\usepackage{hyperref}
\usepackage{xcolor}
\usepackage{colortbl}
\usepackage{array}
\usepackage{enumitem}
\usepackage{needspace}
\usepackage{etoolbox}

\preto\section{\Needspace{8\baselineskip}}

\hypersetup{
  colorlinks=true,
  linkcolor=black,
  urlcolor=blue,
  pdftitle={Modos de conexión TeleCon},
  pdfauthor={TeleCon}
}

\makeatletter
\g@addto@macro\UrlBreaks{\do\_\do\-\do\/\do\:}
\makeatother
\newcommand{\code}[1]{\nolinkurl{#1}}

\definecolor{TblHead}{HTML}{0F6F73}
\definecolor{TblChan}{HTML}{D5F5E3}
\definecolor{TblTelem}{HTML}{D4E6F1}
\definecolor{TblKnob}{HTML}{F5CBA7}

\newcommand{\thc}[1]{\textcolor{white}{\bfseries #1}}
\newcommand{\chip}[2]{\colorbox{#1}{\strut\,\scriptsize #2\,}}
\newcommand{\ui}[1]{\textsf{\textbf{#1}}}

\title{Modos de conexión TeleCon\\[0.3em]
\large Qué elegir antes de conectar}
\author{El ajuste del teléfono, la placa y el sketch deben coincidir}
\date{Versión 1.0 \quad 2026}

\begin{document}
\maketitle

\begin{abstract}
Un modo de conexión es el enlace que el teléfono abre con la ESP32.
Elíjalo en los ajustes de la aplicación antes de pulsar Conectar. El
sketch grabado en la placa debe usar el mismo saludo. Esta nota enumera
cada modo que la aplicación sigue ofreciendo, y la cámara superpuesta,
que no es un modo aparte.
\end{abstract}

{\sffamily\scriptsize
\chip{TblKnob}{Cámara superpuesta}}

\tableofcontents
\vspace{1em}

\section{Elija el modo antes de conectar}

El teléfono no cambia el sketch cuando usted cambia el modo. Solo cambia
cómo habla. Elija el modo y después conecte. Classic Simple, Classic
Binario y BLE Binario están todos disponibles.

\begin{itemize}[leftmargin=1.4cm]
  \item El modo sigue vigente para esa aplicación hasta que lo cambie.
        No hace falta reconectar solo para reasignar un canal de
        telemetría. Sí hace falta reconectar, y volver a grabar, cuando
        cambia el saludo.
\end{itemize}

\section{Los modos de un vistazo}

\begin{center}
\small
\renewcommand{\arraystretch}{1.25}
\setlength{\tabcolsep}{4pt}
\begin{tabular}{>{\raggedright\arraybackslash}p{4.2cm}
                >{\raggedright\arraybackslash}p{2.4cm}
                >{\raggedright\arraybackslash}p{7.2cm}}
\rowcolor{TblHead}
\thc{Modo en la aplicación} & \thc{Placa} & \thc{Sketch} \\
\rowcolor{TblChan}
Classic + Simple & DevKit & \code{TeleCon_Classic_Simple} \\
\rowcolor{TblChan}
Classic + Binario & DevKit & \code{TeleCon_Classic_Binary} \\
\rowcolor{TblChan}
BLE + Binario & DevKit & \code{TeleCon_BLE_Binary} \\
\end{tabular}
\end{center}

La fila opcional \ui{Usar cámara SoftAP} se suma a los tres modos
Bluetooth. No los sustituye. Vea la última sección.

\section{Classic + Simple}

Bluetooth Classic, líneas de texto. Es el primer modo que conviene usar.

\begin{itemize}[leftmargin=1.4cm]
  \item Aplicación: placa \ui{DevKit}, modo
        \ui{Classic + Simple (líneas de texto)}.
  \item Nombre de la placa \code{ESP32-TC-RC-BT-Simple}, PIN \code{1234}.
  \item Saludo \code{RC:CONNECT,proto,simple}.
  \item Palancas, botones y telemetría son líneas como \code{RC:CTRL},
        \code{RC:BTN}, \code{RC:DATA} y \code{RC:PLOT}.
  \item Ocho bytes de telemetría (CH1--CH8) cuando el sketch envía el
        gráfico.
\end{itemize}

\section{Classic + Binario}

La misma radio Classic que Simple, con tramas binarias cortas en lugar de
texto. Más ritmo, los mismos pines, menos que leer en el Monitor Serie.

\begin{itemize}[leftmargin=1.4cm]
  \item Aplicación: placa \ui{DevKit}, modo \ui{Classic + Binario}.
  \item Nombre de la placa \code{ESP32-TC-RC-BT-Binary}.
  \item Saludo \code{RC:CONNECT,proto,binary}.
  \item Las tramas empiezan por \code{AA 55} (palancas), \code{BB 66}
        (botones) y \code{CC 11/22/33} (telemetría).
  \item No empareje esta placa como BLE, y no deje el teléfono en Simple.
\end{itemize}

\section{BLE + Binario}

Bluetooth de baja energía, Nordic UART, con las mismas tramas binarias que
Classic Binario. Úselo en un módulo que solo tenga BLE, como ESP32-C3 o
ESP32-C6, o cuando quiera el radio en reposo más bajo de BLE.

\begin{itemize}[leftmargin=1.4cm]
  \item Aplicación: placa \ui{DevKit}, modo \ui{BLE + Binario}.
  \item Nombre de la placa \code{ESP32-TC-RC-BLE-Binary}.
  \item Saludo \code{RC:CONNECT,proto,binary}.
  \item No hay PIN de Classic. Busque el nombre BLE. No lo empareje como
        un dispositivo Classic.
\end{itemize}

\section{La cámara SoftAP no es un modo}

\ui{Usar cámara SoftAP} es un interruptor sobre un modo Bluetooth de
DevKit. El modo de control sigue siendo Classic Simple, Classic Binario o
BLE Binario.

\begin{itemize}[leftmargin=1.4cm]
  \item Grabe \code{TeleCon_CAM_SoftAP_Video} en una ESP32-CAM aparte.
  \item Únase a \code{TeleCon-RC-CAM-Starter} / \code{telecon1234} solo
        para la imagen.
  \item Las palancas siguen en el modo Bluetooth que ya eligió.
  \item El teléfono no abre el puerto TCP 3333 para esta superposición.
\end{itemize}

La ESP32-CAM no lleva las palancas. Una DevKit hace Bluetooth y una
segunda CAM hace solo el video.

\end{document}
